\section{Introduction}
\label{sec:intro}

Physics-informed neural networks (PINNs) are a widely used deep-learning
surrogate for partial differential equations
\citep{raissi2019physics,karniadakis2021physics,cuomo2022scientific,
lu2021deepxde}. Their appeal is clearest after training: a neural network
represents a mesh-free, differentiable field that can be evaluated at
arbitrary spatial or space-time locations without running a new global
linear or nonlinear solve \citep{leveque2007finite,brenner2008mathematical}.
This amortised evaluation cost is attractive, but it does not by itself
provide the local accuracy guarantees expected of a resolved numerical
discretisation. On several standard PDE benchmarks, PINNs can exhibit large
errors in regimes with multiscale structure, stiffness, sharp gradients, or
unbalanced residual terms \citep{krishnapriyan2021,wang2022whenpinns,
wang2021understanding}. Existing a-posteriori analyses relate residual
norms to global solution errors under stability and approximation
assumptions \citep{mishra2022estimates}; such bounds are valuable, but they
do not directly identify the particular points or subregions where a trained
PINN should be trusted.

The missing object is therefore a \emph{local trust rule}. Empirically, PINN
errors are often highly non-uniform over the domain: a model may approximate
large smooth regions well while failing near fronts, boundary layers,
singular features, or other regions whose solution components are difficult
for the training dynamics to capture \citep{krishnapriyan2021,
wang2022whenpinns,wang2022respecting}. This observation suggests a different
deployment goal from uniform replacement of a classical solver. Instead of
asking a single trained PINN to be accurate everywhere, we ask it to predict
where it is accurate enough to use and where it should abstain. The natural
statistical language for this goal is learning with abstention
\citep{chow1970,bartlett2008classification,Geifman_El-Yaniv_2017,cortes,
cortes2024theory,cao2022generalizing,ni2019calibration,Mao_Mohri_Zhong_2023,
theoretically}: a predictor is paired with a rejector that decides, for each
input, whether to return the prediction or abstain.

Classical abstention would return only a rejection symbol $\bot$ on rejected
inputs. For PDE surrogates this is usually not enough, because downstream
simulation, visualisation, or physical diagnostics often require a complete
solution field. We therefore study abstention with a fallback solver: a
frozen PINN $\widehat{u}_\theta$ supplies the field on accepted points, while
a classical numerical solver is used on rejected regions or queried
locations. If the rejector abstains on a small fraction of the domain, the
hybrid can preserve much of the amortised efficiency of the PINN while using
the solver where the PINN is estimated to be unreliable. The precise cost
depends on the implementation of the fallback solve and the geometry of the
rejected set, so we model it abstractly through an input-dependent
abstention cost $c(x)$ rather than assuming a universal per-point solver
cost.

This perspective is complementary to the large body of work that improves
PINN training. Loss reweighting based on gradient-flow or neural-tangent
kernel diagnostics \citep{wang2021understanding,wang2022whenpinns},
self-adaptive collocation weights \citep{mcclenny2023self}, causal
residual training \citep{wang2022respecting}, residual-based adaptive
sampling \citep{wu2023comprehensive,daw2023mitigating}, and domain
decomposition methods \citep{jagtap2020extended,jagtap2020conservative,
moseley2023finite} can substantially improve particular classes of PDE
problems. However, they still return a single global predictor. They do not,
by themselves, convert a trained PINN into a pointwise certificate of
reliability, nor do they give a post-hoc rule for routing an already-trained
model to a fallback solver. Our goal is to add precisely that missing layer:
a learned rejector whose decision is calibrated to the local error of a
fixed PINN relative to the cost of abstention.

In the setting studied here, the trained PINN is fixed and the learned object
is only the rejector. We write the fixed predictor as
$h=\widehat{u}_\theta$ and measure its pointwise error through a
non-negative discrepancy $\psi(h(x),u^\star(x))$. The ideal rejection rule
compares the conditional risk
$\eta_h(x)=\E[\psi(h(x),Y)\mid X=x]$ with an abstention cost $c(x)$:
accept the PINN when $\eta_h(x)\le c(x)$ and abstain otherwise. This is the
same structural comparison as classical rejection, but in a continuous
regression setting with a spatially varying cost. In practice, $\eta_h$ is
unobserved, so the rejector must be trained from computable local signals
such as the PDE residual and linearised error-transport estimates. The
theoretical task is therefore to design a smooth surrogate whose minimisation
learns the same accept--reject decision as the Bayes rule.

\paragraph{Contributions.}
\begin{itemize}
\item We introduce a post-hoc abstention layer for PINNs: a trained field
  $\widehat{u}_\theta$ is kept frozen, and a separate rejector learns the
  subregions on which the PINN should be replaced by a numerical fallback.
\item We formulate this layer as regression with rejection for continuous
  solution fields, with a general discrepancy $\psi$ and an
  input-dependent abstention cost $c(x)$, rather than the constant-cost
  classification setting usually considered in abstention.
\item We identify the Bayes rejector for a fixed PINN as the local threshold
  rule $\eta_h(x)>c(x)$, where
  $\eta_h(x)=\E[\psi(h(x),Y)\mid X=x]$, making explicit the quantity that a
  PINN reliability model must approximate.
\item We design the rejector around computable physics-based signals,
  combining strong-form residual information with linearised
  error-transport estimates to target the local residual--error gap.
\item We provide the surrogate-consistency and hybrid-solution analysis
  needed for deployment: minimising the smooth rejector loss controls the
  hard abstention decision, and the solver-filled regions yield a complete
  PDE field rather than a rejection symbol.
\end{itemize}
